\documentclass[aspectratio=169,11pt,t]{beamer}

% ============================================================
% Theme and typography
% ============================================================

\usetheme{Madrid}
\usecolortheme{default}
\usefonttheme{professionalfonts}

\setbeamertemplate{navigation symbols}{}
\setbeamertemplate{headline}{}

\setbeamertemplate{footline}{%
  \leavevmode\hbox{%
  \begin{beamercolorbox}[wd=.82\paperwidth,ht=2.5ex,dp=1.0ex,leftskip=.8em]{author in head/foot}
    CSCI 2406 --- Computer Organization \& Architecture (ARM64 / AArch64 Reference)
  \end{beamercolorbox}%
  \begin{beamercolorbox}[wd=.18\paperwidth,ht=2.5ex,dp=1.0ex,rightskip=.8em]{date in head/foot}
    \hfill\insertframenumber/\inserttotalframenumber
  \end{beamercolorbox}}
}

\setbeamersize{
  text margin left=0.65cm,
  text margin right=0.65cm
}

\setbeamerfont{frametitle}{size=\Large, series=\bfseries}
\setbeamerfont{block title}{size=\normalsize, series=\bfseries}
\setbeamertemplate{blocks}[rounded][shadow=false]

% Block vertical padding adjustments
\addtobeamertemplate{block begin}{\vspace{0.1em}}{\vspace{0.1em}}
\addtobeamertemplate{block end}{}{\vspace{0.2em}}

% Base list item separation
\setbeamertemplate{itemize/enumerate body begin}{\setlength{\itemsep}{4pt}}
\setbeamertemplate{itemize/enumerate subbody begin}{\setlength{\itemsep}{2pt}}

% ============================================================
% Packages
% ============================================================

\usepackage[T1]{fontenc}
\usepackage{lmodern}
\usepackage{amsmath}
\usepackage{booktabs}
\usepackage{array}
\usepackage{colortbl}
\usepackage{tikz}

\usetikzlibrary{
  arrows.meta,
  positioning,
  shapes.geometric,
  calc,
  decorations.pathreplacing
}

% ============================================================
% Colors
% ============================================================

\definecolor{navy}{RGB}{24,43,68}
\definecolor{deepblue}{RGB}{45,88,130}
\definecolor{lightblue}{RGB}{235,242,250}
\definecolor{softgray}{RGB}{245,247,250}
\definecolor{darkgray}{RGB}{25,28,32}
\definecolor{gold}{RGB}{200,145,35}

\setbeamercolor{palette primary}{bg=navy,fg=white}
\setbeamercolor{palette secondary}{bg=deepblue,fg=white}
\setbeamercolor{frametitle}{bg=navy,fg=white}
\setbeamercolor{title}{fg=white}
\setbeamercolor{subtitle}{fg=white}
\setbeamercolor{author}{fg=white}
\setbeamercolor{date}{fg=white}
\setbeamercolor{titlelike}{fg=white}
\setbeamercolor{block title}{bg=deepblue,fg=white}
\setbeamercolor{block body}{bg=lightblue,fg=darkgray}
\setbeamercolor{normal text}{fg=darkgray,bg=white}

% ============================================================
% Formatting shortcuts
% ============================================================

\setlength{\parskip}{0pt}

\newcommand{\key}[1]{\textbf{#1}}
\newcommand{\bits}[1]{\texttt{#1}}
\newcommand{\hex}[1]{\texttt{0x#1}}
\newcommand{\code}[1]{\texttt{#1}}

% Section divider macro
\newcommand{\sectiondivider}[2]{%
\begin{frame}[plain]
\begin{tikzpicture}[remember picture,overlay]
\fill[navy] (current page.south west) rectangle (current page.north east);
\fill[gold] (current page.south west) rectangle ([yshift=.15cm]current page.south east);
\end{tikzpicture}
\vspace{1.5cm}
\begin{minipage}{0.85\textwidth}
{\color{gold}\large\bfseries #1\par}
\vspace{0.3cm}
{\color{white}\fontsize{22}{26}\selectfont\bfseries #2\par}
\end{minipage}
\end{frame}
}

% ============================================================
% Metadata
% ============================================================

\title[Topic 6]{
  Data Transfer and Addressing
}

\subtitle{
  Load and store instructions, addressing modes,\\
  memory access alignment, and traversal of arrays and pointers in AArch64
}

\author{
  Computer Organization \& Architecture
}

\date{Topic 6}

% ============================================================
% Presentation Body
% ============================================================

\begin{document}

% --- Slide 1: Title Slide ---
\begin{frame}[plain]
\begin{tikzpicture}[remember picture,overlay]
\fill[navy] (current page.south west) rectangle (current page.north east);
\fill[gold] (current page.south west) rectangle ([yshift=.18cm]current page.south east);
\end{tikzpicture}
\vspace{1.0cm}
\begin{minipage}{0.9\textwidth}
{\color{gold}\large\bfseries CSCI 2406: Topic 06\par}
\vspace{0.3cm}
{\color{white}\fontsize{20}{24}\selectfont\bfseries Data Transfer and Addressing\par}
\vspace{0.3cm}
{\color{lightblue}\normalsize Load/Store architecture, addressing modes, memory alignment rules, and array/pointer traversal in AArch64\par}
\vspace{0.6cm}
{\color{softgray}\small Computer Organization \& Architecture \quad\textbar\quad Topic 6\par}
\end{minipage}
\end{frame}

% --- Slide 2: Topic Outline ---
\begin{frame}{Topic 6 Outline}
    \begin{block}{Topic Objective}
        Analysis of AArch64 load-store memory architecture, memory addressing modes, alignment constraints, and data structure mapping for arrays and pointers.
    \end{block}
    \vspace{0.4em}
    \begin{itemize}
        \item \key{Section 1}: The Load-Store Architecture Paradigm
        \item \key{Section 2}: Load and Store Instructions (\code{ldr}, \code{str})
        \item \key{Section 3}: AArch64 Memory Addressing Modes
        \item \key{Section 4}: Memory Access, Alignment, and Pair Transfers
        \item \key{Section 5}: Traversing Arrays and Pointers
        \item \key{Section 6}: Summary and Worked Self-Test Exercises
    \end{itemize}
\end{frame}

% --- Slide 3: Introductory Concepts ---
\begin{frame}{Bridging Processor Registers and Memory}
    \begin{block}{The Memory Wall Challenge}
        CPU execution speeds outpace main memory bandwidth; explicit data transfer instructions manage movement between high-speed registers and byte-addressed storage.
    \end{block}
    \vspace{0.4em}
    \begin{itemize}
        \item Arithmetic instructions operate exclusively on registers, never directly on memory.
        \item Programs require specialized transfer instructions to read from and write to memory.
        \item Efficient memory access depends on correct addressing mode selection and natural alignment.
    \end{itemize}
\end{frame}

% ============================================================
% SECTION 1
% ============================================================
\sectiondivider{Section 01}{The Load-Store Architecture Paradigm}

% --- Slide 5: Load-Store Architecture Principles ---
\begin{frame}{Load-Store Architecture Principles}
    AArch64 follows a strict RISC load-store design philosophy:
    \vspace{0.2em}
    \begin{itemize}
        \item \key{Separation of Concerns}: Computation instructions (\code{add}, \code{sub}) only process register operands. Memory instructions (\code{ldr}, \code{str}) only move data.
        \item \key{No In-Memory Arithmetic}: An instruction cannot simultaneously fetch a value from memory, add to it, and store it back.
        \item \key{Simplified Datapath}: Keeps hardware control logic streamlined, predictable, and optimized for pipelining.
    \end{itemize}
    \vspace{0.3em}
    \begin{block}{Implication for Compilers}
        C statements like \code{a = b + c;} compile into three distinct steps: load \code{b}, load \code{c}, add them in registers, and store the sum into \code{a}.
    \end{block}
\end{frame}

% --- Slide 6: The Role of Memory Addresses ---
\begin{frame}{Memory Addresses as Register Values}
    Memory locations are targeted using 64-bit virtual addresses held in general-purpose registers:
    \vspace{0.3em}
    \begin{itemize}
        \item The base address of a variable, array, or data structure is computed or loaded into an \code{X} register (e.g., \code{x1}).
        \item Load/store instructions treat this register as a pointer referencing the target byte location.
        \item Arithmetic operations can manipulate pointers dynamically (e.g., advancing through indices).
    \end{itemize}
    \vspace{0.4em}
    \begin{block}{Pointer Register Rule}
        Only 64-bit general-purpose registers (\code{x0}--\code{x30}) can serve as base address pointers in AArch64.
    \end{block}
\end{frame}

% ============================================================
% SECTION 2
% ============================================================
\sectiondivider{Section 02}{Load and Store Instructions}

% --- Slide 8: Basic Load Instruction (LDR) ---
\begin{frame}{Loading Data: The \code{ldr} Instruction}
    The Load Register (\code{ldr}) instruction copies data from memory into a target register:
    \vspace{0.3em}
    \begin{itemize}
        \item \key{Syntax Format}: \code{ldr destination\_reg, [base\_register]}
        \item \key{64-bit Transfer}: \code{ldr x0, [x1]} loads 8 bytes from the address in \code{x1} into \code{x0}.
        \item \key{32-bit Transfer}: \code{ldr w0, [x1]} loads 4 bytes into \code{w0}, zero-extending the upper 32 bits of \code{x0}.
    \end{itemize}
    \vspace{0.4em}
    \begin{block}{Data Sizing Suffixes}
        Sub-word loads are supported via explicit qualifiers like \code{ldrb} (byte) and \code{ldrh} (halfword), with sign-extended variants (\code{ldrsb}, \code{ldrsh}).
    \end{block}
\end{frame}

% --- Slide 9: Basic Store Instruction (STR) ---
\begin{frame}{Storing Data: The \code{str} Instruction}
    The Store Register (\code{str}) instruction copies data from a register into memory:
    \vspace{0.3em}
    \begin{itemize}
        \item \key{Syntax Format}: \code{str source\_reg, [base\_register]}
        \item \key{64-bit Transfer}: \code{str x0, [x1]} writes the 8-byte contents of \code{x0} to the memory address held in \code{x1}.
        \item \key{32-bit Transfer}: \code{str w0, [x1]} writes the lower 4 bytes of \code{w0} to memory.
    \end{itemize}
    \vspace{0.4em}
    \begin{block}{Sub-Word Stores}
        Corresponding store variants \code{strb} (store byte) and \code{strh} (store halfword) allow granular writing of lower register bytes to memory destinations.
    \end{block}
\end{frame}

% --- Slide 10: Size Specifiers and Sign Extension ---
\begin{frame}{Precision Data Sizing and Sign Extension}
    AArch64 load instructions allow precise control over data width and extension behavior:
    \vspace{0.2em}
    \begin{center}
        \small
        \begin{tabular}{p{3.2cm} p{4.2cm} p{4.2cm}}
            \toprule
            \textbf{Instruction} & \textbf{Bytes Loaded} & \textbf{Extension Behavior} \\
            \midrule
            \code{ldrb w0, [x1]} & 1 Byte ($8$ bits) & Zero-extended to 32 bits in \code{w0} \\
            \code{ldrsb x0, [x1]} & 1 Byte ($8$ bits) & Sign-extended to 64 bits in \code{x0} \\
            \code{ldrh w0, [x1]} & 2 Bytes ($16$ bits) & Zero-extended to 32 bits in \code{w0} \\
            \code{ldrsh x0, [x1]} & 2 Bytes ($16$ bits) & Sign-extended to 64 bits in \code{x0} \\
            \bottomrule
        \end{tabular}
    \end{center}
    \vspace{0.2em}
    \begin{block}{Sign Extension Importance}
        When loading signed 8-bit or 16-bit integers into larger registers, using sign-extended variants (\code{ldrsb}/\code{ldrsh}) preserves negative values correctly.
    \end{block}
\end{frame}

% ============================================================
% SECTION 3
% ============================================================
\sectiondivider{Section 03}{AArch64 Memory Addressing Modes}

% --- Slide 12: Overview of Addressing Modes ---
\begin{frame}{Addressing Mode Classification}
    An addressing mode defines how the CPU calculates the target memory address for a load or store instruction:
    \vspace{0.2em}
    \begin{itemize}
        \item \key{Base Register (Indirect)}: Address is taken directly from the base register.
        \item \key{Immediate Offset}: Adds a constant offset to the base register without modifying the base.
        \item \key{Register Offset}: Adds a scaling or offset register to the base address.
        \item \key{Pre-Indexed}: Applies an offset, updates the base register, and accesses memory.
        \item \key{Post-Indexed}: Accesses memory using the base address, then updates the base register.
    \end{itemize}
\end{frame}

% --- Slide 13: Base Register (Indirect) Mode ---
\begin{frame}{Base Register (Indirect) Addressing}
    \begin{block}{Syntax and Semantics}
        $$\code{ldr x0, [x1]}$$
        The memory address is taken strictly from the contents of the base register \code{x1}. The base register remains unchanged after execution.
    \end{block}
    \vspace{0.3em}
    \begin{itemize}
        \item \key{Use Case}: Accessing a variable or dereferencing a pointer whose exact address is already loaded in a register.
        \item Analogous to pointer dereferencing in C: \code{*ptr}.
    \end{itemize}
\end{frame}

% --- Slide 14: Immediate Offset Addressing ---
\begin{frame}{Immediate Offset Addressing}
    \begin{block}{Syntax and Semantics}
        $$\code{ldr x0, [x1, \#16]}$$
        Computes the target address by adding an immediate constant offset (\code{+\#16}) to the base register \code{x1}. The base register \code{x1} is \textbf{not} modified.
    \end{block}
    \vspace{0.3em}
    \begin{itemize}
        \item \key{Use Case}: Accessing fields within a structure or fixed offsets in a local stack frame.
        \item Immediate range rules depend on transfer size (must align with data width constraints).
    \end{itemize}
\end{frame}

% --- Slide 15: Pre-Indexed Addressing Mode ---
\begin{frame}{Pre-Indexed Addressing Mode}
    \begin{block}{Syntax and Semantics}
        $$\code{ldr x0, [x1, \#8]!}$$
        The immediate offset (\code{+\#8}) is added to base register \code{x1}. This new calculated address is used for the memory load, \textbf{and} is written back into \code{x1}.
    \end{block}
    \vspace{0.3em}
    \begin{itemize}
        \item Identified by the trailing exclamation mark (\code{!}).
        \item \key{Use Case}: Updating a pointer before accessing the next sequential memory item.
    \end{itemize}
\end{frame}

% --- Slide 16: Post-Indexed Addressing Mode ---
\begin{frame}{Post-Indexed Addressing Mode}
    \begin{block}{Syntax and Semantics}
        $$\code{ldr x0, [x1], \#8}$$
        Data is loaded from the unmodified address in \code{x1}. After the memory access completes, the offset (\code{+\#8}) is added to \code{x1}, updating the pointer register.
    \end{block}
    \vspace{0.3em}
    \begin{itemize}
        \item Identified by placing the offset outside the square brackets.
        \item \key{Use Case}: Fetching current array elements and automatically advancing the pointer for the next iteration.
    \end{itemize}
\end{frame}

% --- Slide 17: Register Offset Addressing ---
\begin{frame}{Register Offset and Scaled Addressing}
    AArch64 allows an index register to offset the base pointer, optionally scaled by data size:
    \vspace{0.3em}
    \begin{itemize}
        \item \key{Syntax}: \code{ldr x0, [x1, x2]} $\rightarrow$ Address is $\code{x1} + \code{x2}$.
        \item \key{Scaled Register Offset}: \code{ldr x0, [x1, x2, lsl \#3]} $\rightarrow$ Address is $\code{x1} + (\code{x2} \times 8)$.
    \end{itemize}
    \vspace{0.3em}
    \begin{block}{Array Indexing Power}
        Scaling shifts the index register by the exact byte width of the data item ($2^3 = 8$ bytes for doublewords), mapping C array subscripts directly into hardware addressing.
    \end{block}
\end{frame}

% --- Slide 18: Addressing Mode Summary Table ---
\begin{frame}{Addressing Mode Summary}
    \begin{table}[]
        \centering
        \small
        \begin{tabular}{p{2.6cm} p{4.2cm} p{2.4cm} p{2.8cm}}
            \toprule
            \textbf{Mode} & \textbf{AArch64 Assembly} & \textbf{Address} & \textbf{Base Register} \\
            \midrule
            Offset & \code{ldr x0, [x1, x2]} & $\code{x1} + \code{x2}$ & Unchanged \\
            \hline
            Pre-index & \code{ldr x0, [x1, x2]!} & $\code{x1} + \code{x2}$ & $\code{x1} = \code{x1} + \code{x2}$ \\
            \hline
            Post-index & \code{ldr x0, [x1], x2} & $\code{x1}$ & $\code{x1} = \code{x1} + \code{x2}$ \\
            \bottomrule
        \end{tabular}
    \end{table}
    \vspace{0.2em}
    \begin{block}{Note on Register vs. Immediate}
        Offset and index quantities can be specified as general-purpose registers (as shown above) or compatible immediate byte constants (e.g., \code{\#8}).
    \end{block}
\end{frame}

% ============================================================
% SECTION 4
% ============================================================
\sectiondivider{Section 04}{Memory Access, Alignment,\\and Pair Transfers}

% --- Slide 20: Memory Alignment Requirements ---
\begin{frame}{Memory Alignment Rules in AArch64}
    Data alignment dictates how memory addresses map to data sizes:
    \vspace{0.2em}
    \begin{itemize}
        \item \key{Natural Alignment}: An $n$-byte data item must be stored at an address divisible by $n$.
        \begin{itemize}
            \item Doublewords ($8$ bytes) require addresses divisible by $8$.
            \item Words ($4$ bytes) require addresses divisible by $4$.
            \item Halfwords ($2$ bytes) require addresses divisible by $2$.
        \end{itemize}
        \item \key{Strict Enforcement}: Unlike x86, unaligned general-purpose memory accesses in AArch64 can trigger hardware alignment fault exceptions (unless explicitly permitted by specific architectural modes).
    \end{itemize}
\end{frame}

% --- Slide 21: Load and Store Pair Instructions ---
\begin{frame}{Load and Store Pair Instructions (\code{ldp} / \code{stp})}
    AArch64 provides specialized instructions to transfer \textbf{two registers simultaneously}:
    \vspace{0.3em}
    \begin{itemize}
        \item \key{Load Pair (\code{ldp})}: \code{ldp x0, x1, [x2]} loads two 64-bit registers from consecutive memory locations starting at \code{x2}.
        \item \key{Store Pair (\code{stp})}: \code{stp x0, x1, [x2]} stores two 64-bit registers consecutively to memory.
    \end{itemize}
    \vspace{0.3em}
    \begin{block}{Performance Efficiency}
        Pair instructions utilize wide memory data buses, moving 16 bytes in a single instruction cycle and reducing overall instruction count.
    \end{block}
\end{frame}

% --- Slide 22: Stack Operations Using STP and LDP ---
\begin{frame}{Stack Operations via Pair Transfers}
    Stack management relies heavily on pre-index and post-index pair instructions:
    \vspace{0.3em}
    \begin{block}{Pushing to the Stack (Pre-Index)}
        $$\code{stp x29, x30, [sp, \#-16]!}$$
        Decrements stack pointer (\code{sp}) by $16$ bytes, then stores frame pointer (\code{x29}) and link register (\code{x30}).
    \end{block}
    \vspace{0.2em}
    \begin{block}{Popping from the Stack (Post-Index)}
        $$\code{ldp x29, x30, [sp], \#16}$$
        Loads registers from current \code{sp}, then increments stack pointer by $16$ bytes, restoring stack balance.
    \end{block}
\end{frame}

% ============================================================
% SECTION 5
% ============================================================
\sectiondivider{Section 05}{Traversing Arrays and Pointers}

% --- Slide 24: Mapping C Arrays to Assembly ---
\begin{frame}{Mapping C Arrays to AArch64 Assembly}
    Arrays consist of contiguous memory locations holding elements of identical size:
    \vspace{0.2em}
    \begin{itemize}
        \item Let \code{x1} hold the base address of an integer array (\code{int A[]}, where each element is 4 bytes).
        \item Accessing element \code{A[i]} requires computing the effective address: $\text{Base} + (\text{Index} \times \text{Element Size})$.
        \item Assembly implementations utilize scaled register offsets or pointer arithmetic loops.
    \end{itemize}
    \vspace{0.3em}
    \begin{block}{Example Array Access}
        \code{ldr w0, [x1, x2, lsl \#2]} loads \code{A[i]} into \code{w0}, where \code{x2} holds index \code{i} scaled by $2^2 = 4$.
    \end{block}
\end{frame}

% --- Slide 25: Array Traversal via Pointer Arithmetic ---
\begin{frame}{Array Traversal via Pointer Increments}
    Loops often traverse arrays by incrementing pointers directly:
    \vspace{0.3em}
    \begin{itemize}
        \item Initializing base pointer \code{x1} to the start of an array.
        \item Using post-indexed loads to fetch data and advance the pointer in one instruction:
        $$\code{ldr x0, [x1], \#8}$$
        \item Eliminates the overhead of maintaining a separate loop index register.
    \end{itemize}
    \vspace{0.4em}
    \begin{block}{Loop Termination}
        Loops compare the current pointer against an end-of-array sentinel address using comparison (\code{cmp}) and conditional branch instructions.
    \end{block}
\end{frame}

% --- Slide 26: Structs and Record Access ---
\begin{frame}{Accessing Structures and Records}
    Structures group heterogeneous data fields at fixed byte offsets from a base pointer:
    \vspace{0.3em}
    \begin{itemize}
        \item Suppose a struct has fields at offsets $0$ (ID), $4$ (Age), and $8$ (Score).
        \item Base address of struct instance held in pointer register \code{x1}.
        \item Immediate offset addressing accesses individual members without modifying the base pointer:
        \begin{itemize}
            \item \code{ldr w0, [x1, \#0]} \quad (Load ID)
            \item \code{ldr w2, [x1, \#4]} \quad (Load Age)
            \item \code{ldr x3, [x1, \#8]} \quad (Load Score)
        \end{itemize}
    \end{itemize}
\end{frame}

% ============================================================
% SECTION 6
% ============================================================
\sectiondivider{Section 06}{Summary \& Self-Test}

% --- Slide 28: Topic Summary ---
\begin{frame}{Topic 6 Comprehensive Summary}
    \begin{itemize}
        \item \key{Load-Store Paradigm}: Memory access is strictly restricted to explicit \code{ldr} and \code{str} instructions; arithmetic occurs exclusively in registers.
        \item \key{Addressing Modes}: Supports base indirect, immediate offset, register offset, pre-indexed, and post-indexed pointer manipulation.
        \item \key{Alignment Rules}: Natural alignment is strictly enforced to prevent hardware alignment faults.
        \item \key{Pair Transfers}: \code{ldp} and \code{stp} optimize multi-register memory access and stack frame management.
        \item \key{Data Structures}: Scaled offsets and post-indexing enable efficient traversal of arrays and structures.
    \end{itemize}
\end{frame}

% --- Slide 29: Self-Test: Questions 1 \& 2 ---
\begin{frame}{Self-Test: Questions 1 \& 2}
    \begin{block}{Question 1: Load-Store Architecture}
        Why does AArch64 enforce a strict load-store architecture, prohibiting arithmetic instructions from directly operating on memory operands?
    \end{block}
    \vspace{0.4em}
    \begin{block}{Question 2: Immediate vs. Indexed Addressing}
        What is the primary operational difference between immediate offset addressing (\code{[x1, \#8]}) and pre-indexed addressing (\code{[x1, \#8]!})?
    \end{block}
\end{frame}

% --- Slide 30: Self-Test: Solutions 1 \& 2 ---
\begin{frame}{Self-Test: Solutions 1 \& 2}
    \begin{block}{Solution 1}
        Separating load/store operations from computation simplifies hardware datapath design, enables uniform instruction lengths, keeps pipelines predictable, and avoids complex multi-cycle memory-access execution cycles.
    \end{block}
    \vspace{0.3em}
    \begin{block}{Solution 2}
        Immediate offset adds the constant to the base register solely to calculate the target address, leaving the base register unchanged. Pre-indexed addressing calculates the address \textbf{and writes the updated address back} into the base register.
    \end{block}
\end{frame}

% --- Slide 31: Self-Test: Questions 3 \& 4 ---
\begin{frame}{Self-Test: Questions 3 \& 4}
    \begin{block}{Question 3: Post-Indexed Array Traversal}
        Explain how the instruction \code{ldr x0, [x1], \#8} executes step-by-step and why it is useful for array traversal.
    \end{block}
    \vspace{0.4em}
    \begin{block}{Question 4: Memory Alignment}
        What happens when an unaligned 64-bit doubleword load is executed on standard AArch64 hardware?
    \end{block}
\end{frame}

% --- Slide 32: Self-Test: Solutions 3 \& 4 ---
\begin{frame}{Self-Test: Solutions 3 \& 4}
    \begin{block}{Solution 3}
        It loads data from the current address in \code{x1} into \code{x0}, and then increments \code{x1} by $8$ bytes. This automates pointer advancement, making sequential array traversal concise and efficient.
    \end{block}
    \vspace{0.3em}
    \begin{block}{Solution 4}
        Unaligned accesses violate natural alignment rules and typically trigger a hardware alignment fault exception, crashing the program or invoking an OS trap handler.
    \end{block}
\end{frame}

% --- Slide 33: Self-Test: Questions 5 \& 6 ---
\begin{frame}{Self-Test: Questions 5 \& 6}
    \begin{block}{Question 5: Pair Transfers on the Stack}
        How do \code{stp} and \code{ldp} manage stack pointer updates during function entry and exit?
    \end{block}
    \vspace{0.4em}
    \begin{block}{Question 6: Scaled Register Addressing}
        In the instruction \code{ldr x0, [x1, x2, lsl \#3]}, what role does the \code{lsl \#3} modifier play in indexing an array of 64-bit integers?
    \end{block}
\end{frame}

% --- Slide 34: Self-Test: Solutions 5 \& 6 ---
\begin{frame}{Self-Test: Solutions 5 \& 6}
    \begin{block}{Solution 5}
        Using pre-indexed store pair (\code{stp ... [sp, \#-16]!}), the stack pointer is decremented and registers are saved in one atomic step. Post-indexed load pair (\code{ldp ... [sp], \#16}) restores registers and re-increments the stack pointer.
    \end{block}
    \vspace{0.3em}
    \begin{block}{Solution 6}
        The \code{lsl \#3} scales the index in \code{x2} by multiplying it by $2^3 = 8$ (the byte size of a 64-bit doubleword), correctly translating logical array indices into byte offsets.
    \end{block}
\end{frame}

% --- Slide 35: Review: Addressing Mode Summary Table ---
\begin{frame}{Review: Addressing Mode Summary Table}
    \begin{table}[]
        \centering
        \small
        \begin{tabular}{p{3.2cm} p{3.8cm} p{4.5cm}}
            \toprule
            \textbf{Addressing Mode} & \textbf{Assembly Syntax Example} & \textbf{Base Register State After Access} \\
            \midrule
            \key{Base Register} & \code{ldr x0, [x1]} & Unchanged \\
            \key{Immediate Offset} & \code{ldr x0, [x1, \#16]} & Unchanged \\
            \key{Pre-Indexed} & \code{ldr x0, [x1, \#8]!} & Updated ($\text{Base} + 8$) \\
            \key{Post-Indexed} & \code{ldr x0, [x1], \#8} & Updated ($\text{Base} + 8$) \\
            \bottomrule
        \end{tabular}
    \end{table}
\end{frame}

% --- Slide 36: Wrap-Up \& Transition to Topic 7 ---
\begin{frame}{Wrap-Up \& Transition to Topic 7}
    \begin{block}{Moving Forward to Topic 7}
        Now that we have mastered memory data transfer, addressing modes, and alignment, our next session will explore:
        \begin{itemize}
            \item \key{Control Flow \& Branching}: Unconditional and conditional branches in AArch64.
            \item \key{Subroutine Linkage}: Function calls (\code{bl}), return linkage (\code{ret}), and ABI parameter passing conventions.
        \end{itemize}
    \end{block}
    \vspace{0.3em}
    \begin{center}
        \textbf{End of Topic 6 --- Questions \& Discussion}
    \end{center}
\end{frame}

% --- Slide 37: Additional Source Notes ---
\begin{frame}{Topic 6 Source Notes}
    \begin{block}{Source Reference Verification}
        \begin{itemize}
            \item Load-store mechanics and offset calculations follow standard ARMv8/ARMv9 AArch64 ISA architecture guidelines.
            \item Stack frame management conforms to standard Procedure Call Standard for the ARM 64-bit Architecture (AAPCS64).
        \end{itemize}
    \end{block}
\end{frame}

\end{document}